\section*{Chapitre \ref{ch:rs:live-experiments} --- Expérimentations en production}

\begin{solution}{exo:rs:live-experiments:1}
$-0.55 \pm 1.96 \times 0.48$~: de $-1.49$ à 0.39. L'intervalle contient zéro, et il contient aussi des économies trois fois supérieures à toute valeur plausible~: l'expérience est
non concluante, et non négative. Il aurait fallu plus d'ordres, une part plus grande, ou \omterm{def:rs:live-experiments:cuped}{CUPED}.
\end{solution}

\begin{solution}{exo:rs:live-experiments:2}
$z = (2\,485 - 2\,400)/\sqrt{24\,000 \times 0.1 \times 0.9} = 85/46.5 = 1.83$, une p-value bilatérale de 0.07~: pas de déséquilibre à 5~\%. À une p-value de 0.0001, il faut cesser d'interpréter
l'estimation et chercher la cause (ordres perdus, dupliqués ou réacheminés dans un bras, filtre appliqué après l'assignation, serveur exécutant une ancienne règle d'assignation)~;
la comparaison n'est pas valable tant que la répartition n'est pas expliquée.
\end{solution}

\begin{solution}{exo:rs:live-experiments:3}
$2.80 \times 0.085 = 0.24$ point de base~; la puissance pour un effet de 0.3 est de 0.94.
\end{solution}

\begin{solution}{exo:rs:live-experiments:4}
Autorisés~: l'estimation avant transaction (fixée avant l'assignation), le mouvement du marché sur la fenêtre prévue (hors de portée de l'algorithme), le
shortfall de la veille du titre (avant l'assignation). Non autorisés~: la durée réalisée et le taux d'exécution, que l'algorithme modifie~; les ajuster
revient à retirer une partie de l'effet. Avec $R^2 = 0.51$, la variance est multipliée par 0.49~: $77 \times 0.49 = 38$ jours.
\end{solution}

\begin{solution}{exo:rs:live-experiments:5}
$n = 3^2 \times 23.1^2/(5^2(1-s))$~: 194 ordres et 8.1 jours à 1~\%, 214 ordres et 0.89 jour (348 minutes sur une séance de 6.5 heures) à 10~\%, 385 ordres et
0.32 jour à 50~\%. Le défaut de Knight envoyait des ordres enfants sans égard pour la quantité exécutée de l'ordre parent~: un invariant (la quantité d'un ordre enfant ne dépasse jamais celle de son parent)
vérifié à chaque ordre, une limite de position calibrée sur l'étape, et un kill switch auraient arrêté le défaut en quelques secondes. Un ordre hors de contrôle est une règle violée, et non
un petit décalage de la moyenne.
\end{solution}

\begin{solution}{exo:rs:live-experiments:6}
$-0.30 + 0.25 \times 0.883 = -0.079$, soit environ $-0.08$. Randomisés ordre par ordre, les ordres traités et les ordres témoins ont la même proportion attendue de
voisins traités, si bien que le débordement est le même dans les deux bras et s'annule dans la différence. L'unité titre-jour (tous les ordres d'un même titre un jour donné, dans un même bras) contient
l'\omterm{def:rs:live-experiments:ab}{interférence}, et mesure le \omterm{def:rs:live-experiments:canary}{déploiement progressif}.
\end{solution}

\begin{solution}{exo:rs:live-experiments:7}
Médiane du nombre de jours avant arrêt~: 56 pour $\tau = 0.1$, 38 pour $\tau = 0.3$, 44 pour $\tau = 1$. Un petit $\tau$ place la masse de l'alternative sur des effets trop faibles pour
être distingués rapidement~; une valeur élevée la répartit sur des effets bien supérieurs à 0.3 et paie pour eux un facteur $\sqrt{V/(V+\tau^2)}$. Le meilleur $\tau$ est proche de
l'effet que l'on attend.
\end{solution}

\begin{solution}{exo:rs:live-experiments:8}
Des consultations quotidiennes d'un test à horizon fixe, sans aucun effet, atteignent $p < 0.05$ à une consultation quelconque dans les 23 jours, dans 25~\% des expériences et non 5~\%. Il faut soit fixer
l'horizon et ne regarder qu'une fois, soit utiliser une p-value toujours valide (le \omterm{def:rs:live-experiments:mde}{test séquentiel} par mélange), dont le niveau tient à chaque consultation.
\end{solution}

\begin{solution}{pb:rs:live-experiments:1}
\begin{enumerate}
\item Moyenne 11.7, écart-type 23.1 points de base~: l'estimation avant transaction, son erreur, le mouvement signé du marché sur la fenêtre, le bruit idiosyncratique et un
effet de journée.
\item Déterministe, sans état, identique sur chaque serveur, indépendant d'une expérience à l'autre, reproductible a posteriori~: aucune table partagée ne peut se désynchroniser.
\item 2\,485 contre 2\,400 prévus~: $z = 1.83$, p-value 0.07, pas de déséquilibre.
\item L'économie directe, $-0.30$~; le \omterm{def:rs:live-experiments:canary}{déploiement progressif} économise $-0.08$.
\item Que chaque serveur exécutait le nouveau code (une seconde relecture), que les alertes parvenaient à quelqu'un, et que les limites et invariants de chaque étape bornaient ce qu'elle pouvait faire.
\item 77 jours (la formule donne 77.7).
\item L'estimation avant transaction et le mouvement du marché sur la fenêtre prévue~; $R^2 = 0.51$~; 38 jours.
\item 72 et 70 jours~: l'estimation avant transaction n'explique que 9~\% de la variance.
\item Un spread de 0.66 et 3\,253 jours pour 0.3 point de base~: une unité par jour, et l'effet de journée est le bruit.
\item Spread de 0.093 et 35 jours pour 0.3 point de base, comme le meilleur plan au niveau de l'ordre~; il estime l'effet de $-0.08$ du \omterm{def:rs:live-experiments:canary}{déploiement progressif}, dont la détection prendrait 503 jours.
\item 29~\% et 33~\%.
\item 46 et 62 sur 4\,000.
\item La moitié d'ici le jour 38, 4~\% d'ici le jour 10, 98~\% d'ici le jour 120.
\item Pour une médiane de 6 jours.
\item \textbf{Résultat à retenir.} 77 jours de séance à la moitié du flux sans \omterm{def:rs:live-experiments:cuped}{CUPED}, 38 avec (la formule donne 77.7 et 38.2).
\item Sur ses propres ordres, oui~: $-0.30$ avec une erreur type de 0.085 en 60 jours avec \omterm{def:rs:live-experiments:cuped}{CUPED}. Pour la maison, à peine~: tout basculer économise 0.08, ce qui n'est pas
distinguable de zéro en deux années d'activité.
\item Par un \omterm{def:rs:live-experiments:canary}{déploiement canari} avec invariants et limites à l'échelle, puis une expérience randomisée par titre-jour et analysée au sein de la journée avec \omterm{def:rs:live-experiments:cuped}{CUPED}, une règle séquentielle,
des garde-fous sur le taux d'exécution et les mark-outs, une procédure de retour arrière à chaque étape~; et se demander si une économie de 0.08 vaut le changement.
\item 194 ordres et 8.1 jours~: la statistique sert aux petits décalages~; s'appuyer sur des invariants par ordre, des limites dimensionnées pour le \omterm{def:rs:live-experiments:canary}{déploiement canari} et un kill switch pour les défauts de grande ampleur.
\item La métrique et les garde-fous, l'unité, la part, les covariables, l'horizon ou la règle séquentielle et $\tau$, l'\omterm{def:rs:live-experiments:mde}{effet minimal détectable}, les décisions à
chaque issue, consignées dans le \omterm{def:rs:the-research-process:log}{journal de recherche}.
\item La différence que le changement produit sur les \omterm{def:rs:live-experiments:ab}{unités de randomisation} sous l'assignation effectivement appliquée, qui égale l'effet de tout basculer que si
le résultat d'aucune unité ne dépend de l'assignation d'une autre.
\end{enumerate}
\end{solution}

\begin{solution}{iq:rs:live-experiments:1}
\omterm{def:rs:live-experiments:paper}{Trading simulé} pour la parité (signaux, contrôles de risque, calendrier complet)~; un \omterm{def:rs:live-experiments:canary}{déploiement canari} vérifié, des invariants par ordre, des limites ajustées à sa part et un
kill switch~; un \omterm{def:rs:live-experiments:ab}{test A/B} dont l'\omterm{def:rs:live-experiments:ab}{unité de randomisation} intègre une \omterm{def:rs:live-experiments:ab}{interférence}, \omterm{def:rs:live-experiments:cuped}{CUPED} et une règle séquentielle~; une montée en charge par étapes avec retour arrière testé. La perte de Knight
en 2012 provenait d'un \omterm{def:rs:live-experiments:canary}{déploiement progressif} dont le contrôle de l'étape (un serveur sur huit) avait échoué et dont les avertissements n'étaient pas des alertes.
\end{solution}

\begin{solution}{iq:rs:live-experiments:2}
Un ajustement par régression sur des covariables non affectées par le traitement~: la variance diminue de $1 - R^2$ (Bing indiquait environ la moitié~; ici 0.51 avec le mouvement du marché).
Elle échoue si une covariable est affectée par le traitement (durée, taux d'exécution~: biais), ou si elle explique peu la variance (l'estimation avant transaction seule~: 9~\%).
\end{solution}

\begin{solution}{iq:rs:live-experiments:3}
Chaque consultation offre une nouvelle chance de faux positif~: des consultations quotidiennes donnent 29~\% de faux positifs en 38 jours. Il faut fixer l'horizon, ou utiliser des p-values toujours valides issues d'un
\omterm{def:rs:live-experiments:mde}{test séquentiel} du rapport de vraisemblance par mélange, qui borne l'erreur sur l'ensemble des consultations (inégalité de Ville).
\end{solution}

\begin{solution}{iq:rs:live-experiments:4}
\omterm{def:rs:live-experiments:ab}{Interférence}~: le gain du nouvel algorithme est venu en partie au détriment des autres ordres de la maison sur les mêmes titres (il a négocié plus tôt et leur a laissé davantage d'impact).
La randomisation au niveau de l'ordre annule le débordement entre les bras~; le \omterm{def:rs:live-experiments:canary}{déploiement progressif} le subit. Randomiser par titre-jour (ou par client, ou par jour) pour mesurer le \omterm{def:rs:live-experiments:canary}{déploiement progressif}.
\end{solution}

\begin{solution}{iq:rs:live-experiments:5}
Hacher un identifiant stable avec le nom de l'expérience (SHA-256, les premiers octets ramenés à une fraction comparée à la part)~: le même ordre reçoit le même bras sur chaque serveur
sans état partagé, et les expériences sont indépendantes. Vérifier la répartition avec le test de \omterm{def:rs:live-experiments:ab}{déséquilibre du ratio d'échantillonnage} avant de lire un résultat, et enquêter sur un
déséquilibre avant de faire confiance à la comparaison.
\end{solution}

\begin{solution}{iq:rs:live-experiments:6}
$n = ((z_{1-\alpha/2} + z_{1-\beta})\sigma/\Delta)^2 (1/s + 1/(1-s))$~; à 5~\% et 80~\%, le facteur vaut 2.80. Avec $\sigma = 23.1$, $\Delta = 0.3$ et une
répartition égale~: $n = (2.80 \times 23.1/0.3)^2 \times 4 \approx 186\,500$ ordres, 77.7 jours sur 2\,400.
\end{solution}
